\section{Proofs and extensions}
\label{app:proofs}
\label{sec:proofs}

The results below make the observation assumptions explicit and provide
verifiable foundations for the implementation. They are applications of
established identification, convex duality, concentration, and decision-theoretic
arguments, not claims of new general mathematical principles. Relevant
antecedents include partial monitoring \citep{kirschner2023}, partial
identification under measurement error \citep{finkelstein2021}, optimal
recovery \citep{donoho1994}, and comparison of experiments
\citep{blackwell1953}.

Let $A\in\mathbb R^{m\times K}$ be nonnegative and column-stochastic. Write
$\mathcal S_K=\{p\in\mathbb R^K:p\geq0,\ \mathbf1^\top p=1\}$,
$q=Ap$, and $\Delta(p)=d^\top p$ for a fixed known vector $d\in\mathbb R^K$.
For $q\in A\mathcal S_K$, define
\[
 \mathcal P_q=\{p\in\mathcal S_K:Ap=q\},\qquad
 L(q)=\min_{p\in\mathcal P_q}d^\top p,\quad
 U(q)=\max_{p\in\mathcal P_q}d^\top p.
\]
These extrema are attained because $\mathcal P_q$ is nonempty and compact.
The target is a population contrast. It is not recovery of individual private
preferences.

\begin{theorem}[Global and local identification]
\label{thm:global-identification}
The contrast $d^\top p$ is identified for every feasible $q$ if and only if
$d\in\operatorname{row}(A)$. At a particular $q$, it is identified if and only if
$d\perp\operatorname{span}(\mathcal P_q-\mathcal P_q)$.
\end{theorem}
\begin{proof}
If $d=A^\top v$, then $d^\top p=v^\top q$ for every compatible population.
Conversely, if $d\notin\operatorname{row}(A)$, there is an $h\in\ker(A)$ with
$d^\top h\ne0$. Column normalization gives $\mathbf1^\top h=0$.
Choose an interior $p_0\in\mathcal S_K$ and sufficiently small $t>0$ so that
$p_0\pm th\geq0$. Both populations belong to the simplex, produce $Ap_0$, and
have different contrasts. The local statement follows because constancy of
$d^\top p$ on $\mathcal P_q$ is equivalent to orthogonality to every pairwise
difference, hence to their span.
\end{proof}

The global condition need not hold at an identified boundary point. For
example, the channel with columns $(1,0)^\top,(0,1)^\top,(0,1)^\top$ cannot
generally separate classes two and three, but $q=(1,0)^\top$ uniquely determines
$p=(1,0,0)^\top$. This distinction prevents a rank-only test from replacing the
feasible-set calculation. Observation-span conditions have direct antecedents
in partial monitoring \citep{kirschner2023}.

\begin{theorem}[Magnitude of the globally hidden contrast]
\label{thm:hidden-width}
Define
\[
 W(A,d)=\max_{p,r\in\mathcal S_K:Ap=Ar}|d^\top(p-r)|.
\]
Then
\begin{equation}
\label{eq:hidden-width-dual}
 W(A,d)
 =\max_{Ah=0,\ \|h\|_1\leq2}d^\top h
 =2\min_{v\in\mathbb R^m}\|d-A^\top v\|_\infty.
\end{equation}
\end{theorem}
\begin{proof}
Every $h=p-r$ in the first optimization satisfies $Ah=0$ and $\|h\|_1\leq2$.
Conversely, $Ah=0$ implies $\mathbf1^\top h=0$. Let $h_+,h_-$ be its positive
and negative parts, with common mass $s=\|h\|_1/2\leq1$. For any
$u\in\mathcal S_K$, the vectors $p=h_++(1-s)u$ and $r=h_-+(1-s)u$ are
compatible simplex points and satisfy $p-r=h$. The feasible set is symmetric,
so maximizing the absolute contrast equals maximizing the signed contrast.

For the second equality, consider
\[
 \min_{v,t}\ t\quad\text{subject to}\quad
 d-A^\top v\leq t\mathbf1,\quad
 A^\top v-d\leq t\mathbf1,\quad t\geq0.
\]
Assign nonnegative multipliers $a,b$ to these two inequalities. Minimizing the
Lagrangian over unrestricted $v$ requires $A(a-b)=0$; minimizing over $t\geq0$
requires $\mathbf1^\top(a+b)\leq1$. The dual objective is $d^\top(a-b)$.
Writing $z=a-b$ gives precisely $Az=0,\ \|z\|_1\leq1$: the reverse direction
uses $a=z_+,b=z_-$. Primal feasibility follows by taking $v=0$ and
$t=\|d\|_\infty$; the primal is bounded below, so finite-dimensional LP strong
duality applies. Rescaling $h=2z$ proves the identity.
\end{proof}

$W(A,d)$ is a worst-case width over all possible $q$; the interval at an observed
$q$ may be narrower. The dual quantity is a standard distance to an observable
subspace, with conceptual antecedents in optimal recovery
\citep{donoho1994}. It is not, by itself, a measured market effect.

\begin{theorem}[Irreducible estimation and decision loss]
\label{thm:decision-loss}
Fix $q$ and write $L=L(q)$, $U=U(q)$, and $w=U-L$. Suppose the data consist of
any finite number of iid actions from $q$, together with independent analyst
randomization. Every estimator has worst-case mean absolute error at least
$w/2$ and worst-case mean squared error at least $w^2/4$ over $\mathcal P_q$.
If $q$ is supplied exactly, the midpoint attains both bounds.

If $L<0<U$, the minimax customer-value regret for choosing between variants zero
and one, when $q$ is known exactly, is
\begin{equation}
\label{eq:minimax-regret}
 R^*(q)=\frac{-LU}{U-L}.
\end{equation}
It is attained by choosing variant one with probability $t^*=U/(U-L)$.
The minimax regret among deterministic decisions is $\min\{-L,U\}$. If the
interval does not strictly straddle zero, minimax regret is zero.
\end{theorem}
\begin{proof}
All compatible populations yield the same data distribution. Let $T$ have the
common distribution of an estimator at populations attaining the endpoints.
Pointwise,
\[
 |T-L|+|T-U|\geq w,\qquad
 \frac{(T-L)^2+(T-U)^2}{2}
 =\left(T-\frac{L+U}{2}\right)^2+\frac{w^2}{4}.
\]
Taking expectations bounds the larger endpoint risk from below. The constant
midpoint attains both lower bounds over the whole interval when $q$ is known.

For a randomized decision rule let $t$ be its common probability of selecting
variant one. Its regrets at the lower and upper endpoints are $(-L)t$ and
$U(1-t)$, respectively; intermediate contrasts give no larger regret. Minimizing
$\max\{(-L)t,U(1-t)\}$ over $t\in[0,1]$ equates the terms and gives
\eqref{eq:minimax-regret}. Restricting $t$ to zero or one gives the deterministic
result. If the interval lies weakly on one side of zero, a single variant is
weakly optimal throughout.
\end{proof}

For strict opposite-sign endpoints, the decision error probabilities are $t$
and $1-t$, so their equally weighted average is $1/2$. These are fixed-channel
indistinguishability statements, a standard lower-bound argument
\citep{lattimore2020}. The attainability claims presume exact $q$; finite-sample
uncertainty can add loss. A probe that changes the observation channel can
break the indistinguishability.

\begin{theorem}[Coarsening]
\label{thm:coarsening}
Let $G$ be column-stochastic and set $B=GA$. For every feasible $q$,
\[
 \mathcal P(A,q)\subseteq\mathcal P(B,Gq).
\]
The interval under $B$ at $Gq$ therefore contains the interval under $A$ at $q$,
and $W(B,d)\geq W(A,d)$.
\end{theorem}
\begin{proof}
$Ap=q$ implies $Bp=GAp=Gq$. Minimization over the enlarged set cannot increase
the lower bound, and maximization cannot decrease the upper bound. Every pair
with $Ap=Ar$ also has $Bp=Br$, proving the global inequality.
\end{proof}

This finite-channel implication belongs to the comparison-of-experiments
viewpoint \citep{blackwell1953}. It does not establish that a more competent
agent is a coarsening of another agent. That is an additional, testable
relationship, not an assumption licensed by higher task success.

\begin{theorem}[A bias-aware finite-sample certificate]
\label{thm:linear-certificate}
Let $Y_1,\ldots,Y_n$ be iid categorical actions from $q=Ap$. Choose
$v\in\mathbb R^m$ independently of these data and define
$\epsilon(v)=\|d-A^\top v\|_\infty$ and
$\operatorname{rng}(v)=\max_a v_a-\min_a v_a$.
With probability at least $1-\alpha$,
\begin{equation}
\label{eq:linear-certificate}
 \left|\frac1n\sum_{i=1}^n v_{Y_i}-d^\top p\right|
 \leq \epsilon(v)+\operatorname{rng}(v)
       \sqrt{\frac{\log(2/\alpha)}{2n}}.
\end{equation}
If the deployed channel is $A'$ and
$\max_k\operatorname{TV}(A'_{\cdot k},A_{\cdot k})\leq\tau$, the bound holds
for iid data from $A'p$ after adding $\tau\operatorname{rng}(v)$.
\end{theorem}
\begin{proof}
The mean of $v_{Y_i}$ under $A$ is $v^\top Ap$. Since $p$ is a probability
vector, its difference from $d^\top p$ has absolute value at most
$\|A^\top v-d\|_\infty$. Hoeffding's inequality for variables in
$[\min_a v_a,\max_a v_a]$ bounds the sampling deviation by the second term
with probability at least $1-\alpha$. The triangle inequality proves
\eqref{eq:linear-certificate}.

For distributions $r,s$, write $b=\min_a v_a$ and use
$\sum_a(r_a-s_a)=0$ to obtain
$|v^\top(r-s)|\leq\operatorname{rng}(v)\operatorname{TV}(r,s)$.
Applying this bound to every channel column and averaging over $p$ bounds the
additional drift bias by $\tau\operatorname{rng}(v)$.
\end{proof}

Minimizing the right-hand side over $v$ using only $A,d,n,\alpha,\tau$ preserves
the guarantee. Choosing $v$ against the same field observations requires a
separate argument. An estimated channel is not automatically a known channel;
its uncertainty must be covered by a valid drift bound or the joint confidence
set. Bias-aware linear estimation and concentration are established tools
\citep{donoho1994,lattimore2020}.

\begin{theorem}[Identification from logged probes]
\label{thm:logged-probes}
Let probe $e\in\{1,\ldots,E\}$ be assigned independently of latent class with
probability $\rho_e>0$, where $\sum_e\rho_e=1$, and let its channel be $A_e$.
Suppose the probe identity is recorded and all probes share the same population
composition $p$. The joint probe/action channel is
$B=\operatorname{vstack}(\rho_1A_1,\ldots,\rho_EA_E)$, and
\[
 d^\top p\text{ is globally identified}
 \quad\Longleftrightarrow\quad
 d\in\operatorname{row}(B)
 =\operatorname{span}\!\left(\bigcup_e\operatorname{row}(A_e)\right).
\]
\end{theorem}
\begin{proof}
$B$ is column-stochastic and
$\ker(B)=\bigcap_e\ker(A_e)$ because every $\rho_e$ is positive. Apply
Theorem~\ref{thm:global-identification} and orthogonal-complement identities.
\end{proof}

If probe labels are discarded, the channel is instead $\sum_e\rho_e A_e$.
For example, $A_1=I_2$ and $A_2$ equal to $I_2$ with its rows exchanged are each
fully revealing. At equal assignment probabilities their unlogged average has
identical columns $(1/2,1/2)^\top$. Logged experimental variation can therefore
be informative when unlogged variation is not. Probe choice and its costs
connect directly to experimental design and partial monitoring
\citep{kirschner2023,lattimore2020}.

\begin{proposition}[Continuous cost allocation for fixed weights]
\label{prop:cost-allocation}
Suppose independent samples from probe $e$ have size $n_e>0$, and fixed weights
$v_e$ satisfy $\sum_e A_e^\top v_e=d$. Write
$\sigma_e^2=\operatorname{Var}_{A_ep}(v_{e,Y})$ and assume $\sigma_e>0$.
The estimator $\sum_e n_e^{-1}\sum_i v_{e,Y_{ei}}$ is unbiased, with variance
$\sum_e\sigma_e^2/n_e$. Under positive costs $c_e$ and the continuous budget
$\sum_e c_en_e=C$, its minimum variance and optimal allocation are
\[
 V_{\min}=\frac{(\sum_e\sigma_e\sqrt{c_e})^2}{C},\qquad
 n_e^*=\frac{C\sigma_e/\sqrt{c_e}}{\sum_j\sigma_j\sqrt{c_j}}.
\]
\end{proposition}
\begin{proof}
The weight identity gives unbiasedness. Independence gives the variance.
Cauchy--Schwarz yields
\[
 (\sum_e\sigma_e\sqrt{c_e})^2
 \leq (\sum_e\sigma_e^2/n_e)(\sum_e c_en_e).
\]
Equality holds at the displayed allocation.
\end{proof}
This familiar calculation is not a new acquisition algorithm. Integer sample
sizes, unknown variances, adaptive weights, or required minimum allocations
need additional treatment. Zero variances require the corresponding limiting
allocation or explicit minimum-count constraints.

\begin{proposition}[Rectangular uncertainty and coverage]
\label{prop:uncertainty-coverage}
Suppose channel columns lie in coordinate boxes
$\underline A_{ak}\leq A_{ak}\leq\overline A_{ak}$, with $0\leq\underline A\leq\overline A\leq1$ and each containing at least
one simplex vector. Frequencies lie in
$\underline q\leq q\leq\overline q$. Introduce $J_{ak}$ and impose
\begin{align}
 p&\geq0,\quad J\geq0, & \mathbf1^\top p&=1, \nonumber\\
 \sum_a J_{ak}&=p_k, &
 \underline A_{ak}p_k\leq J_{ak}&\leq\overline A_{ak}p_k,\label{eq:joint-lp}\\
 \underline q_a&\leq\sum_kJ_{ak}\leq\overline q_a. &&\nonumber
\end{align}
Minimizing and maximizing $d^\top p$ over these constraints gives attainable
bounds within the supplied rectangular model. If the channel and frequency
boxes cover their true values with probabilities at least $1-\alpha_A$ and
$1-\alpha_q$, respectively, the interval covers $d^\top p$ with probability
at least $1-\alpha_A-\alpha_q$.
\end{proposition}
\begin{proof}
Every admissible $(A,p)$ maps to $J=A\operatorname{diag}(p)$ satisfying
\eqref{eq:joint-lp}. Conversely, if $p_k>0$, define $A_{ak}=J_{ak}/p_k$.
The resulting column sums to one and satisfies its box constraints. If
$p_k=0$, the constraints force the entire column of $J$ to zero; choose any
simplex vector from that column's nonempty channel box. These choices produce
an admissible channel with frequencies $\sum_kJ_{\cdot k}$. This establishes
exactness, including zero-mass columns.

On the intersection of the two coverage events, the true pair
$(p,A\operatorname{diag}(p))$ is feasible. Its contrast lies between the extrema.
The union bound gives the asserted probability without requiring independence
of the confidence events.
\end{proof}

This confidence-set propagation is already present in the measurement-error
literature, including the supplementary Proposition~1 of
\citet{finkelstein2021}. An infeasible program is a failure state, not a narrow
confidence interval. Fixed-sample intervals do not justify optional stopping.

\paragraph{Uncertain outcomes.}
If $d$ lies in an independent coordinate rectangle
$[\underline d,\overline d]$, sharp rectangular-model endpoints are
$\min_p\underline d^\top p$ and $\max_p\overline d^\top p$ over the same
feasible set. This follows because $p\geq0$ makes the coordinate endpoints
optimal for each fixed $p$. If this outcome rectangle has error probability
$\alpha_d$, the same proof gives coverage at least
$1-\alpha_A-\alpha_q-\alpha_d$. Parameter correlations excluded by the
rectangle can make these bounds conservative.

\paragraph{Omitted classes.}
A separate contamination model is needed if some customers lie outside the
calibrated taxonomy. Let their mass be $\eta\in[0,\bar\eta]$, their total action
masses be $u_a\geq0$, and known-class masses be $z_k\geq0$. Replace the mass and
frequency constraints by
\[
 \sum_k z_k=1-\eta,\quad \sum_a u_a=\eta,\quad
 \underline q_a\leq\sum_kJ_{ak}+u_a\leq\overline q_a,
\]
and replace $p$ by $z$ in the channel constraints. If unknown-class contrasts
lie in $[-D,D]$, introduce $t$ with $-D\eta\leq t\leq D\eta$ and optimize
$d^\top z+t$. These constraints are linear. Every allowed contaminated model
maps to these masses. Conversely, normalize nonzero known-class columns as in
the proof above; when $\eta>0$, one unknown class with channel $u/\eta$ and
contrast $t/\eta$ realizes the residual. At $\eta=0$, both residuals vanish.
Thus the extension is exact for this broad contamination model. Merely widening
an interval computed from contaminated frequencies without modifying the
observation constraints does not establish such protection.

\begin{proposition}[Calibration on a fixed heterogeneous panel]
\label{prop:heterogeneous-panel}
For class $k$, let $Y_{k1},\ldots,Y_{kn_k}$ be independent responses to a
prespecified, possibly heterogeneous prompt panel. Define
\[
 \bar A_{ak}=\frac1{n_k}\sum_i\Pr(Y_{ki}=a),\qquad
 \widehat A_{ak}=\frac1{n_k}\sum_i\mathbf1\{Y_{ki}=a\}.
\]
For $m$ categories and $K$ classes, simultaneous coordinate intervals with
radius
\[
 r_k=\sqrt{\frac{\log(2mK/\alpha_A)}{2n_k}}
\]
cover $\bar A$ with probability at least $1-\alpha_A$. Endpoints may be clipped
to $[0,1]$.
\end{proposition}
\begin{proof}
For fixed $a,k$, the indicators are independent variables in $[0,1]$, though
they need not have identical expectations. Hoeffding's inequality gives
\[
 \Pr(|\widehat A_{ak}-\bar A_{ak}|>r_k)
 \leq 2e^{-2n_kr_k^2}=\frac{\alpha_A}{mK}.
\]
A union bound over all coordinates proves simultaneous coverage. Clipping to
the probability range cannot exclude a true coordinate already covered.
\end{proof}

The result targets the fixed-panel average channel, not an arbitrary deployment
channel. Matching prompt composition or an explicit transport bound is still
required. Alternatively, independently sampled prompts from a prespecified
distribution, combined with independent stable responses, support the iid
multinomial interpretation used by the benchmark's calibration intervals.
Refusals and errors must remain declared categories or enter a missing-data
model. Dependence across calls or adaptive prompt selection is not covered by
the elementary panel argument. Simulated profiles remain synthetic experimental
inputs, irrespective of model fluency.
